% sections/00_executive_summary.tex
% 执行摘要（不编号）

\section*{执行摘要}
\addcontentsline{toc}{section}{执行摘要}

\subsection*{核心结论}

本报告从物理、工程、商业三个递进层面，对低地球轨道（LEO）AI 计算数据中心的可行性进行独立技术分析。核心结论如下：

\begin{itemize}[leftmargin=*]
    \item \textbf{物理可行，但已逼近当前材料天花板。} 斯特藩-玻尔兹曼定律不否定 1 MW 级轨道散热——AI1 卫星所宣称的 1,400\;W/m\textsuperscript{2}（双面辐射、$\varepsilon\approx0.90$、$T\approx342.5$\;K）在物理上自洽。然而，该工作点几乎用尽了当前高辐射率热控涂层在寿命末期（EOL）退化后的可用裕度。散热密度翻倍在 2035 年前无单一技术路径可支撑。（完整推导见第4章、附录 C\;\S C.1）

    \item \textbf{工程脆弱，单星质量与可靠性是主要瓶颈。} 将计算硬件送入太空并维持其在轨运行，是一个多约束耦合的工程问题。计算载荷质量链已从此前行业分析中使用的"千瓦每吨"反推法（该方法因分母来源不可靠和定义域歧义已被永久废弃，详见附录 E）重构为以 NVIDIA GB300 NVL72 地面硬件为锚点的逐项累加法。基准场景单星总质量约 8.48\;t，其中光伏阵列（55.3\%）和系统裕量/冗余（21.9\%）合计占制造成本的绝大部分。（完整推导见第5章、附录 C\;\S C.3）

    \item \textbf{商业以场景分层，不存在"单一正确数字"。} 本报告建立了三个可比经济分支：10年-GaAs 主链基线总成本约\$0.77B（区间\$0.55B--\$1.1B，取决于关键假设变动）（单代 1\;MW 持续 IT 负载、10 年总窗口、\$200/kg 发射单价）；5年-GaAs 对照分支因需整星重射一次，总成本升至约 \$1,341.26M；5年-HJT 并行路线补齐估计约 \;\$1,005M，但受 AI1 70\;m 翼展几何约束排除，无法作为可落地主路线。发射成本在 TCO 中占比仅约 1.8\%，真正决定商业可行性的因素是光伏阵列成本、在轨寿命和硬件可靠性。（完整结果见第6章及附录 D 的脚本路由）
\end{itemize}

\subsection*{主链 TCO 数字（速查）}

\begin{table}[H]
\centering
\caption{三分支 TCO 总览（共同分母：1\;MW 持续 IT 负载、10 年总窗口）}
\label{tab:exec_tco}
\footnotesize
\begin{tabular}{lc>{\centering\arraybackslash}p{2.8cm}>{\centering\arraybackslash}p{4.4cm}}
\toprule
\textbf{分支} & \textbf{单星质量} & \textbf{10年总成本} & \textbf{状态} \\
\midrule
10年-GaAs（主链） & 8.48\;t & \textbf{\$764.98M}\tablefootnote{按SpaceX公开120kW持续IT口径为星数分母、LEO轨道电源链闭合后的10年-GaAs模型输出；应理解为"官方宣称口径下的乐观可比基线"，非对未来项目成本的审计值。} & 当前唯一可完整给出的单代主链 \\
5年-GaAs（对照）  & 7.10\;t & \$1,341.26M & 寿命缩短至 5 年，需重射一次 \\
5年-HJT（并行）   & 9.13\;t & $\approx$\;\$1,005M\quad（补齐估计） & 70\;m翼展几何排除；数值为补齐估计，非下界 \\
\bottomrule
\end{tabular}
\end{table}

总成本公式：$\Sigma$（制造成本 $+$ 发射成本 $+$ 保险（卫星制造价值$\times$3--5\% + 发射成本$\times$8\%；当前模型仅计入发射成本部分作为下界估计）$+$ 运维（\$5M/年$\times 10$年）$+$ 退役（\$3M/代$\times$代数））。单代卫星当量 $= \text{1\;MW} / \text{120\;kW} \approx 8.33$。详见第6章相关表格及附录 D\;\S \ref{sec:script-entry}。

\subsection*{分歧定性}

NVIDIA CEO 黄仁勋与 SpaceX CEO 埃隆·马斯克关于太空算力的公开立场差异，构成了本报告的核心研究动机。

但两人的分歧\emph{不是"做不做"}。两人均认同太空算力在文明尺度上具有长期必然性；真正的分歧在于\emph{"多快能做"和"瓶颈在哪里"}：

\begin{itemize}[leftmargin=*]
    \item 黄仁勋从芯片工程视角出发，强调散热通量、宇宙射线、太阳风暴和微陨石等多参数耦合的工程难度，以"NVIDIA 已在太空部署 CUDA"证明自己并非反对太空算力，而是对进度有更保守的判断。
    \item 马斯克从文明能源视角出发，将轨道 AI 嵌套进卡尔达肖夫文明等级叙事，隐含假设所有相关参数（发射成本、光伏效率、电池寿命、制造规模）会同步改善。
    \item 本报告的独立判断：\textbf{同步改善是可能的，但多参数同步改善的概率不足以支撑近期明确承诺。} 分歧本质不是"谁算错了"，而是双方使用了不同的时间窗、功率口径、架构代际和系统边界——当这些差异被参数化后，分歧从"方向争议"降维为"参数假设差异"。（口径对齐矩阵见附录 A）
\end{itemize}

\subsection*{阅读指南}

本报告采用模块化结构，不同读者可按需选读：

\begin{itemize}[leftmargin=*]
    \item \textbf{仅需结论的读者}：读完本章即可。本章给出的三层判断和主链 TCO 数字是全文的核心产出。
    \item \textbf{需了解背景的读者}：继续阅读第1章（太空算力的定义、关键参与方和 2024--2026 年事件时间线）。
    \item \textbf{需理解分析框架的读者}：阅读第1--2章。第2章定义了全文的五大统一口径（功率、时间窗、架构代际、成本边界、发射单价）和三层递进框架，是理解后续各章的前提。
    \item \textbf{需核实事实基础的读者}：阅读第3章（12 条一手言论的完整汇总与分歧界定）和附录 D（五层证据路由表，可从正文结论追溯到原始 URL）。
    \item \textbf{需审查技术判断的读者}：阅读第4--6章和附录 A（口径对齐矩阵）、附录 B（47 参数逐条裁决）、附录 C（核心计算方法与公式）。
    \item \textbf{需了解完整判断链的读者}：通读全文及全部五篇附录。
\end{itemize}

\subsection*{章节与附录索引}

\begin{itemize}[leftmargin=*]
    \item \textbf{第1章 背景与研究对象}：太空算力的定义、提出动机、关键参与方（SpaceX/NVIDIA）、2024--2026 年时间线、研究范围与边界。
    \item \textbf{第2章 分析方法与比较基线}：物理-工程-商业三层递进框架、五大统一口径定义、口径对齐矩阵初版。
    \item \textbf{第3章 一手言论与分歧界定}：马斯克 7 条言论汇总表、黄仁勋 5 条言论汇总表、分歧定性（假设差异而非对错）。
    \item \textbf{第4章 科学可行性}：斯特藩-玻尔兹曼定律推导与散热面积、LEO 真实热环境约束、散热密度边界。
    \item \textbf{第5章 工程可行性}：bottom-up 计算载荷质量链（替代 kW/t 反推链）、GaAs/HJT 电源路线对照、电池三分支、整星质量拆解、在轨寿命与维修、发射运力现实检验。
    \item \textbf{第6章 商业可行性}：三分支 TCO 完整拆解与对比、口径差异的量化影响、翻转条件分析。
    \item \textbf{第7章 综合判断}：已知事实 / 可计算链 / 当前不可判定的三层收束、时间线展望、结论稳健性说明。
    \item \textbf{附录 A}：口径对齐矩阵（7 个场景 $\times$ 9 个维度的完整比较）。
    \item \textbf{附录 B}：47 参数逐条裁决表（档位、取值、裁决状态、来源路由）。
    \item \textbf{附录 C}：核心计算方法与公式（S-B 散热推导、轨道电源能量预算、bottom-up 质量链）。
    \item \textbf{附录 D}：五层证据路由表（结论 $\to$ 交叉核验 $\to$ JSON $\to$ 裁决记录 $\to$ URL/公式）。
    \item \textbf{附录 E}：已废弃推导链说明（kW/t 反推链废弃原因、禁线清单、参数使用边界声明）。
\end{itemize}

\subsection*{证据体系说明}

本报告涉及的全部 47 个关键参数已按四档证据体系逐条分级并裁决：

\begin{itemize}[leftmargin=*]
    \item \textbf{一档（物理推导链，5 参数）}：由明确物理公式和边界条件直接推导，如太阳常数、散热面积。
    \item \textbf{二档（权威资料，7 参数）}：来自可回查的权威公开来源，如 AI1 功率参数（多源交叉验证）、AzurSpace 4G32 GaAs 效率。
    \item \textbf{三档（可迁移证据，17 参数）}：来自可迁移的行业资料，如 Starlink 电池 DoD/比能量、地面 HJT 组件批发价。
    \item \textbf{四档（合理推断，14 参数）}：公开资料不足、基于工程推断，如散热器面密度、平台系数、PCDU 比功率等——涉及约 35\% 的制造成本项。四档参数占比 30\%（14/47）是 AI1 信息公开不足的客观约束，而非分析缺陷。
\end{itemize}

正文中每个关键数字附证据档位标注；完整参数裁决表见附录 B；追溯路由见附录 D。

\subsection*{方法论声明}

本报告从物理定律（斯特藩-玻尔兹曼）出发，经工程硬件锚点（NVIDIA GB300 NVL72），到独立 Python 建模的完整计算链，遵循第一性原理自底向上验证的方法论。TCO 中光伏阵列（$\sim$50\%）已通过二档/三档交叉锚定；系统裕量/冗余/附加（AIT裕量）仍是主导成本项（$\sim$20\%），但其取值仍带有四档推断属性。即使四档参数的取值发生合理变动，方向性结论不会发生数量级翻转。本报告追求的不是预测精度，而是结论的稳健性。

\endinput
